\section{Construction of the Original POD Basis}
\label{sec:original_pod_basis}

This section describes the construction of the original proper
orthogonal decomposition (POD) basis from the full-order temperature
solutions. The procedure consists of generating a parameterized
training dataset, assembling the snapshot matrix, performing singular
value decomposition, and selecting a reduced set of spatial modes.
The resulting basis is subsequently used in the Galerkin
reduced-order model.

\subsection{Original training dataset}

The original training dataset is generated using a uniform
two-dimensional parameter grid over the prescribed ranges of the
effective parameter $\mu$ and the Brinkman number $Br$. The
parameter grid contains ten equally spaced values for each
parameter:
\[
N_{\mu}=10,\qquad N_{Br}=10.
\]
The parameter ranges are
\[
1\leq\mu\leq100,\qquad
0.01\leq Br\leq2.
\]
The Cartesian product of the two parameter grids gives a total of
\[
\boxed{N_s=N_{\mu}N_{Br}=100}
\]
parameter combinations.

For each parameter pair $(\mu_k,Br_k)$, the full-order solver
computes a steady-state temperature field $\mathbf{T}_k$. A
simulation case is accepted into the training dataset only when
the solver reports successful convergence. The simulation summary
stores the case identifier, parameter values, convergence status,
final residual, number of iterations, and computation time.
All 100 cases in the original training dataset are marked as
successful.

The resulting dataset therefore contains 100 full-order
temperature solutions associated with 100 distinct parameter
combinations. An earlier 25-case parameter sweep was also
performed during the development of the project; however, the
POD basis and energy results presented in this section correspond
to the subsequent 100-case dataset.

\subsection{Snapshot preparation}

The POD basis is constructed using temperature deviations rather
than the complete temperature fields. As introduced in the
previous section, the temperature field is decomposed as
\[
T_k(x,y)=T_{\mathrm{ref}}(y)+\theta_k(x,y),
\qquad
T_{\mathrm{ref}}(y)=y.
\]
The reference temperature satisfies the prescribed
non-homogeneous temperature boundary conditions. Subtracting
this fixed profile gives the deviation field
\[
\boxed{
\boldsymbol{\theta}_k
=
\mathbf{T}_k-\mathbf{T}_{\mathrm{ref}}.
}
\]
This transformation separates the fixed reference profile
from the parameter-dependent part of the temperature field.
The resulting deviation fields satisfy homogeneous temperature
boundary conditions and are used to construct the snapshot
matrix.

The snapshot preparation procedure consists of the following
steps:
\begin{enumerate}
    \item The converged physical temperature field is loaded
    for each accepted simulation case.
    \item Ghost cells are removed to retain only the physical
    computational domain.
    \item The reference temperature field is subtracted from
    the physical temperature field.
    \item The resulting two-dimensional deviation field is
    flattened into a one-dimensional vector of length
    $N_xN_y$.
    \item The vector is stored as one column of the snapshot
    matrix.
\end{enumerate}

The flattening operation follows the same indexing convention
as the full-order solver, with the wall-normal index varying
fastest. The snapshot matrix is consequently defined as
\[
\mathbf{X}
=
\begin{bmatrix}
| & | & & |\\
\boldsymbol{\theta}_1 &
\boldsymbol{\theta}_2 &
\cdots &
\boldsymbol{\theta}_{100}\\
| & | & & |
\end{bmatrix}.
\]

For the present computational grid,
\[
N_x=50,\qquad N_y=20,
\]
and therefore each snapshot contains 1000 physical degrees of
freedom. Since the training dataset contains 100 snapshots,
the resulting snapshot matrix has dimensions
\[
\boxed{
\mathbf{X}\in\mathbb{R}^{1000\times100}.
}
\]

The parameter values are stored separately from the snapshot
matrix as metadata. Their association with the snapshots can
be represented by the parameter matrix
\[
\mathbf{P}
=
\begin{bmatrix}
\mu_1 & \mu_2 & \cdots & \mu_{100}\\
Br_1 & Br_2 & \cdots & Br_{100}
\end{bmatrix}
\in\mathbb{R}^{2\times100}.
\]
Thus, column $k$ of $\mathbf{X}$ corresponds to column $k$
of $\mathbf{P}$:
\[
\boldsymbol{\theta}_k
\longleftrightarrow
\begin{bmatrix}
\mu_k\\
Br_k
\end{bmatrix}.
\]
This correspondence is maintained through the simulation
metadata and case identifiers. Although the parameter values
are not directly used in the standard POD decomposition,
they are required to associate each snapshot with its
corresponding parameter combination.

\subsection{Singular value decomposition}

The snapshot matrix is decomposed using the economy singular
value decomposition (SVD):
\[
\boxed{
\mathbf{X}=\mathbf{U}\mathbf{\Sigma}\mathbf{W}^{T}.
}
\]
For the present snapshot matrix, the dimensions of the
decomposition are
\[
\begin{aligned}
\mathbf{X}&\in\mathbb{R}^{1000\times100},\\
\mathbf{U}&\in\mathbb{R}^{1000\times100},\\
\mathbf{\Sigma}&\in\mathbb{R}^{100\times100},\\
\mathbf{W}^{T}&\in\mathbb{R}^{100\times100}.
\end{aligned}
\]

The matrices obtained from the decomposition have the
following roles:
\begin{itemize}
    \item $\mathbf{U}$ contains the spatial POD modes as its
    columns. Each column represents a spatial pattern in the
    temperature-deviation field.

    \item $\mathbf{\Sigma}$ is a diagonal matrix containing
    the singular values, ordered from largest to smallest:
    \[
    \sigma_1\geq\sigma_2\geq\cdots\geq\sigma_{100}\geq0.
    \]
    The magnitude of each singular value indicates the
    contribution of its corresponding mode to the snapshot
    matrix.

    \item $\mathbf{W}$ contains the right singular vectors.
    These vectors describe the contributions of the individual
    snapshots to the decomposition and provide the coefficients
    used to reconstruct the original snapshots.
\end{itemize}

A reduced spatial basis is obtained by retaining only the first
$r$ columns of $\mathbf{U}$:
\[
\boxed{
\mathbf{V}_r=\mathbf{U}(:,1:r).
}
\]
Equivalently,
\[
\mathbf{V}_r
=
\begin{bmatrix}
\boldsymbol{\phi}_1 &
\boldsymbol{\phi}_2 &
\cdots &
\boldsymbol{\phi}_r
\end{bmatrix},
\]
where $\boldsymbol{\phi}_i$ denotes the $i$-th spatial POD
mode.

The columns of $\mathbf{U}$ are orthonormal. Therefore, the
retained modes satisfy
\[
\boxed{
\mathbf{V}_r^T\mathbf{V}_r=\mathbf{I}_r.
}
\]
This orthogonality property allows the temperature-deviation
field to be represented by a reduced set of independent
spatial basis functions.

\subsection{POD energy and mode selection}

The singular values quantify the relative contribution of the
spatial modes to the snapshot matrix. The energy fraction
associated with the $i$-th mode is defined as
\[
\boxed{
E_i=
\frac{\sigma_i^2}
{\displaystyle\sum_{j=1}^{N_s}\sigma_j^2}.
}
\]
The cumulative energy captured by the first $r$ modes is
\[
\boxed{
E_r=
\frac{\displaystyle\sum_{i=1}^{r}\sigma_i^2}
{\displaystyle\sum_{j=1}^{N_s}\sigma_j^2}.
}
\]
The squared singular values are used because their sum
represents the squared Frobenius norm of the snapshot matrix.
Consequently, $E_r$ measures the fraction of the total
snapshot energy represented by the first $r$ modes.

The singular values and corresponding energy fractions
obtained from the original 100-snapshot dataset are reported
in Table~\ref{tab:pod_energy}. The first eight modes are
included to show the decay of the singular-value spectrum.

\begin{table}[htbp]
    \centering
    \caption{Singular values and POD energy distribution
    of the original 100-snapshot dataset.}
    \label{tab:pod_energy}
    \begin{tabular}{cccc}
        \toprule
        \textbf{Mode} &
        \textbf{Singular value} &
        \textbf{Energy fraction} &
        \textbf{Cumulative energy} \\
        \midrule
        1 & 788.0597 & 0.926360 & 0.926360 \\
        2 & 215.0846 & 0.069005 & 0.995365 \\
        3 & 54.73393 & 0.004469 & 0.999833 \\
        4 & 10.41370 & 0.000162 & 0.999995 \\
        5 & 1.810038 & $4.89\times10^{-6}$ & $\approx 1$ \\
        6 & 0.247505 & $9.14\times10^{-8}$ & $\approx 1$ \\
        7 & 0.039928 & $2.38\times10^{-9}$ & $\approx 1$ \\
        8 & 0.004025 & $2.42\times10^{-11}$ & $\approx 1$ \\
        \bottomrule
    \end{tabular}
\end{table}

The first mode accounts for 92.6360\% of the total snapshot
energy. Retaining two modes increases the cumulative energy
to 99.5365\%, while the first three modes capture 99.9833\%.
With four retained modes, the cumulative energy reaches
99.9995\%. The cumulative values for modes five to eight
are displayed approximately as unity because of rounding.

The rapid decay of the singular values indicates that the
original training snapshots can be represented with high
energy retention using a relatively small number of spatial
modes. In particular, the first three or four modes provide
a highly compact representation of the training dataset.

However, cumulative POD energy alone does not establish
the accuracy of the reduced-order model at unseen parameter
combinations. It quantifies how well the retained modes
represent the original snapshot matrix, rather than how
accurately the reduced model predicts solutions outside the
training dataset. Prediction accuracy must be evaluated
separately using the reduced-order model.

\subsection{POD reconstruction versus new-parameter prediction}

The SVD provides the modal coefficients required to
reconstruct the snapshots used to construct the POD basis.
To obtain these coefficients, define the truncated singular
value matrix and right singular-vector matrix as
\[
\mathbf{\Sigma}_r
=
\operatorname{diag}(\sigma_1,\ldots,\sigma_r),
\qquad
\mathbf{W}_r=\mathbf{W}(:,1:r).
\]

For an existing snapshot $k$, its reduced coefficient vector
is obtained directly from the truncated SVD:
\[
\boxed{
\mathbf{a}_k
=
\mathbf{\Sigma}_r\mathbf{W}_r^T(:,k).
}
\]
The corresponding temperature-deviation field is
reconstructed through a linear combination of the retained
spatial modes:
\[
\hat{\boldsymbol{\theta}}_k
=
\mathbf{V}_r\mathbf{a}_k.
\]
Adding the reference temperature profile gives the
reconstructed full temperature field:
\[
\boxed{
\hat{\mathbf{T}}_k
=
\mathbf{T}_{\mathrm{ref}}
+
\mathbf{V}_r\mathbf{a}_k.
}
\]
This procedure is referred to as POD reconstruction of an
existing snapshot. The coefficient vector is directly
available from the SVD because the snapshot is already a
column of the original snapshot matrix.

The situation is different for a new parameter combination
$(\mu^*,Br^*)$ that was not included in the training dataset.
There is no corresponding snapshot column in $\mathbf{X}$
and, consequently, the SVD does not provide its modal
coefficient vector. The spatial basis alone is therefore
insufficient to determine the temperature field at this new
parameter combination.

An additional relationship between the parameters and the
reduced coefficients is required:
\[
\boxed{
(\mu,Br)\longrightarrow\mathbf{a}(\mu,Br).
}
\]
In the present work, this relationship is obtained by
projecting the discretized governing energy equation onto
the retained POD basis. The resulting reduced system
determines the modal coefficients for a specified parameter
combination. Once these coefficients have been calculated,
the corresponding temperature field is reconstructed as
\[
\boxed{
\hat{\mathbf{T}}(\mu,Br)
=
\mathbf{T}_{\mathrm{ref}}
+
\mathbf{V}_r\mathbf{a}(\mu,Br).
}
\]
This distinction separates the direct reconstruction of
existing training snapshots from the prediction of
temperature fields at new parameter combinations. The
following section develops the Galerkin projection and
formulates the resulting reduced-order model.

\subsection{Code implementation}

The construction and analysis of the original POD basis
are implemented through a sequence of Python scripts.
The workflow is divided into parameter-grid generation,
full-order simulations, snapshot preparation, SVD,
reconstruction, and visualization.

\begin{itemize}
    \item \texttt{F01\_parameter\_grid.py}: Defines the
    parameter limits and the number of sampling points for
    each parameter. The uniform parameter grids are generated
    using \texttt{numpy.linspace} and combined using
    \texttt{numpy.meshgrid} and
    \texttt{numpy.column\_stack}. The resulting parameter
    combinations are saved to a CSV file.

    \item \texttt{F02\_parameter\_sweep.py}: Loads the
    parameter-grid CSV file and iterates over the parameter
    combinations. For each case, it reads $\mu$ and $Br$,
    calls \texttt{solve\_with\_direct} to obtain the
    steady-state temperature field, and records selected
    simulation information. The summary, including the
    parameter values and convergence information, is saved
    as a CSV file. This stage generates the full-order
    dataset used for POD.

    \item \texttt{F03\_pod\_analysis/f01\_snapshot\_matrix.py}:
    Loads the simulation summary and selects the converged
    cases. For each case, it loads the corresponding
    temperature field, subtracts the reference temperature,
    removes ghost cells, and flattens the resulting
    two-dimensional deviation field into a vector. These
    vectors are stacked column-wise to construct the
    snapshot matrix, which is saved as a NumPy array.

    \item \texttt{F03\_pod\_analysis/f02\_pod.py}: Loads the
    snapshot matrix and performs its SVD using
    \texttt{numpy.linalg.svd}. It calculates the modal and
    cumulative energy distributions from the singular
    values and saves the resulting POD data, including
    the spatial modes and singular values.

    \item \texttt{F03\_pod\_analysis/f03\_visualize\_pod\_modes.py}:
    Loads the spatial modes from the saved POD results,
    determines the number of available modes, and generates
    plots of the individual spatial modes.

    \item \texttt{F03\_pod\_analysis/f04\_pod\_reconstruction.py}:
    Loads the POD results and snapshot metadata, then
    truncates the SVD according to the selected number of
    retained modes. It calculates the reconstructed snapshot
    matrix and evaluates the reconstruction error for each
    training case. The reconstructed temperature fields
    and error records are saved together with their
    corresponding parameter information.

    \item \texttt{F03\_pod\_analysis/f05\_visualize\_reconstruction.py}:
    Loads the simulation summary and the full-order and
    reconstructed temperature fields for three selected
    cases. It generates comparative plots of the
    full-order temperature fields, POD reconstructions,
    and their corresponding error fields.
\end{itemize}

This workflow separates the generation of the full-order
training dataset from the construction and assessment of
the POD basis. The reconstruction scripts assess how
accurately the retained modes reproduce the snapshots
used to construct the basis. The prediction of solutions
at new parameter combinations is subsequently handled by
the Galerkin reduced-order model.